% from main.tex
%%%%%%%%%%  scan idx 411  |  folio 398  %%%%%%%%%%
\text{avec } 0\le\alpha_i\le n-i.
\]
De plus l'additivité de \(\lambda_t\) entraine que \(\prod_i(1+\xi_i t)=\lambda_t(E)\), et les relations du corollaire sont vérifiées. Si \(C\) est la \(K^{\bullet}(S)\)-algèbre engendrée par des générateurs \(\xi_i\), soumis aux relations~:
\[
  \sigma_i(\xi_1,\ldots,\xi_n)=\lambda^i(E),
\]
il existe donc un homomorphisme surjectif de \(C\) sur \(K^{\bullet}(X)\). D'après 4.4, cet homomorphisme transforme une base de \(C\) en une base de \(K^{\bullet}(X)\), c'est donc un isomorphisme, ce qui achève la démonstration du corollaire.

Traitons le cas général d'un fibré en drapeaux de type \(\pi\). Le fibré en drapeaux usuel \(D(E)\) peut s'obtenir à partir de \(D_\pi(E)\) en construisant successivement les fibrés en drapeaux des \(F_j\) (4.2). Il en résulte que \(K^{\bullet}(D(E))\), considéré comme module sur \(K^{\bullet}(X)\), admet une base formée des éléments considérés dans 4.3.i), les \(\xi_i^{(j)}\) étant les classes des faisceaux analogues aux \(F_j\) intervenant dans le scindage des images inverses des \(F_j\) sur \(D(E)\) (notations de 3.5 et 3.7). Donc \(K^{\bullet}(X)\) est un sous-anneau de \(K^{\bullet}(D(E))\) ; de plus, l'additivité de \(\lambda_t\) entraine que~:
\[
  \lambda_i^{(j)}=\sigma_i(\xi_1^{(j)},\ldots,\xi_{p_j}^{(j)}),
\]
et \(K^{\bullet}(X)\) contient la sous-\(K^{\bullet}(S)\)-algèbre de \(K^{\bullet}(D(E))\) engendrée par les \(\lambda_i^{(j)}\). Soit \(B\) cette sous-algèbre ; d'après 4.3.i), et ce qui vient d'être dit, on a dans \(K^{\bullet}(D(E))\) une famille d'éléments qui forme une base à la fois pour la structure de \(B\)-module et pour la structure de \(K^{\bullet}(X)\)-module ; par suite, \(K^{\bullet}(X)=B\), ce qui montre que \(K^{\bullet}(X)\) est engendré par les \(\lambda_i^{(j)}\). L'idéal des relations entre les \(\lambda_i^{(j)}\) est alors donné par 4.3.ii), compte
